\paragrafos{Ευθεία και κύκλος}{Β.3.13, «Σχετικές θέσεις ευθείας και κύκλου»}

\begin{anagke}
Στο κέντρο του κήπου η θεία Μαρία έχει στρογγυλό ποτιστικό που ρίχνει νερό σε απόσταση $4$ μέτρων προς όλες τις κατευθύνσεις. Κοντά στο ποτιστικό περνά ένα ευθύγραμμο μονοπάτι και η Ζωή ρωτά: «Θα βραχούμε περπατώντας;». Ο Ηλίας μετρά την απόσταση του μονοπατιού από το ποτιστικό: $2{,}5$ μέτρα. Ο Βασίλης λέει ότι θα βραχεί σε δύο σημεία, ο Νίκος ότι σε ένα και η Φανή ότι θα προσέξει μήπως δεν βραχεί καθόλου. «Εξαρτάται», λέει η κυρία Αλεξάνδρα, «από το πόσο μακριά περνά το μονοπάτι». Η Ρήνη βρίσκει ότι υπάρχει μια ειδική απόσταση όπου το μονοπάτι μόλις και μετά βίας ακουμπά το νερό.

Σε αυτή την παράγραφο βρίσκουμε πόσα κοινά σημεία μπορεί να έχουν μια ευθεία και ένας κύκλος και πώς σχετίζεται αυτό με την απόσταση της ευθείας από το κέντρο.
\end{anagke}

\begin{lexi}{τέμνουσα, εφαπτομένη, εξωτερική}{θέσεις ευθείας και κύκλου}
Η \lexiword{τέμνουσα} «κόβει» τον κύκλο. Η \lexiword{εφαπτομένη} είναι αυτή που «εφάπτεται», δηλαδή απλώς αγγίζει (το ρήμα «άπτομαι» σημαίνει «αγγίζω»). Η \lexiword{εξωτερική} ευθεία μένει «έξω» από τον κύκλο.
\end{lexi}

\begin{metafrasi}
Η ευθεία και ο κύκλος έχουν δύο κοινά σημεία. & τέμνουσα, $\mathrm{ΟΜ}<\rho$\\
Η ευθεία και ο κύκλος έχουν ένα κοινό σημείο. & εφαπτομένη, $\mathrm{ΟΜ}=\rho$\\
Η ευθεία και ο κύκλος δεν έχουν κοινό σημείο. & εξωτερική, $\mathrm{ΟΜ}>\rho$\\
\end{metafrasi}

\tmima{Η θεωρία}

Ένας κύκλος $(\mathrm{Ο},\rho)$ και μια ευθεία $\varepsilon$ μπορεί να έχουν \textbf{δύο}, \textbf{ένα} ή \textbf{κανένα} κοινό σημείο. Ονομάζουμε $\mathrm{ΟΜ}$ την \textbf{απόσταση} του κέντρου $\mathrm{Ο}$ από την ευθεία $\varepsilon$, δηλαδή το κάθετο τμήμα από το $\mathrm{Ο}$ προς την $\varepsilon$.

\begin{center}
\begin{tikzpicture}[scale=0.85]
\draw[thick] (0,0) circle (1.3);
\fill (0,0) circle (1.6pt) node[above left] {\small $\mathrm{Ο}$};
\draw[thick,gblue] (-2,-0.7)--(2,-0.7); \draw[dashed] (0,0)--(0,-0.7);
\fill (0,-0.7) circle (1.4pt) node[below right] {\small $\mathrm{Μ}$};
\node at (0,-2.9) {\small τέμνουσα, $\mathrm{ΟΜ}<\rho$};
\begin{scope}[xshift=4.8cm]
\draw[thick] (0,0) circle (1.3);
\fill (0,0) circle (1.6pt) node[above left] {\small $\mathrm{Ο}$};
\draw[thick,gblue] (-2,-1.3)--(2,-1.3); \draw[dashed] (0,0)--(0,-1.3);
\fill (0,-1.3) circle (1.4pt) node[below right] {\small $\mathrm{Μ}$};
\node at (0,-2.9) {\small εφαπτομένη, $\mathrm{ΟΜ}=\rho$};
\end{scope}
\begin{scope}[xshift=9.6cm]
\draw[thick] (0,0) circle (1.3);
\fill (0,0) circle (1.6pt) node[above left] {\small $\mathrm{Ο}$};
\draw[thick,gblue] (-2,-1.9)--(2,-1.9); \draw[dashed] (0,0)--(0,-1.9);
\fill (0,-1.9) circle (1.4pt) node[below right] {\small $\mathrm{Μ}$};
\node at (0,-2.9) {\small εξωτερική, $\mathrm{ΟΜ}>\rho$};
\end{scope}
\end{tikzpicture}
\end{center}

\begin{kanonas}{σχετικές θέσεις ευθείας και κύκλου}
Αν $d=\mathrm{ΟΜ}$ είναι η απόσταση του κέντρου από την ευθεία $\varepsilon$, τότε:
\begin{itemize}[leftmargin=1.2em,itemsep=0.3ex,topsep=2pt]
\item αν $d<\rho$, η $\varepsilon$ είναι \textbf{τέμνουσα} του κύκλου και έχουν \textbf{δύο} κοινά σημεία,
\item αν $d=\rho$, η $\varepsilon$ είναι \textbf{εφαπτομένη} του κύκλου και έχουν \textbf{ένα} κοινό σημείο, το \textbf{σημείο επαφής},
\item αν $d>\rho$, η $\varepsilon$ είναι \textbf{εξωτερική} και \textbf{δεν έχουν} κοινό σημείο.
\end{itemize}
\end{kanonas}

Όταν η $\varepsilon$ εφάπτεται στο σημείο $\mathrm{Μ}$, το $\mathrm{ΟΜ}$ είναι ακτίνα και η απόστασή της από το $\mathrm{Ο}$, άρα:

\begin{idiotita}{εφαπτομένη}
Η εφαπτομένη ενός κύκλου σε ένα σημείο του είναι \textbf{κάθετη στην ακτίνα} που καταλήγει στο σημείο επαφής. Από κάθε σημείο του κύκλου περνά \textbf{μία μόνο} εφαπτομένη.
\end{idiotita}

Αν από ένα σημείο $\mathrm{Μ}$ έξω από τον κύκλο φέρουμε τις δύο εφαπτόμενες που ακουμπούν στα $\mathrm{Α}$ και $\mathrm{Β}$, τα τμήματα $\mathrm{ΜΑ}$ και $\mathrm{ΜΒ}$ λέγονται \textbf{εφαπτόμενα τμήματα}.

\begin{prosoxi}
\begin{itemize}[leftmargin=1.2em,itemsep=0.3ex,topsep=2pt]
\item Η \emph{απόσταση} από το κέντρο είναι το κάθετο τμήμα $\mathrm{ΟΜ}$, όχι οποιοδήποτε τμήμα προς την ευθεία.
\item Όσο μικρότερη είναι η απόσταση, τόσο περισσότερο «μπαίνει» η ευθεία στον κύκλο: μικρή απόσταση σημαίνει τέμνουσα, όχι εξωτερική.
\item Η ευθεία που περνά από το κέντρο είναι τέμνουσα (απόσταση $0$) και όχι εφαπτομένη.
\item Η εφαπτομένη είναι κάθετη στην \emph{ακτίνα} του σημείου επαφής, όχι σε οποιαδήποτε ακτίνα.
\end{itemize}
\end{prosoxi}

\tmima{Λυμένα παραδείγματα}

\paradeigma{Η θέση της ευθείας}
\tixero Κύκλος $(\mathrm{Ο},5)$ cm και τρεις ευθείες με αποστάσεις από το $\mathrm{Ο}$: $3$ cm, $5$ cm και $7$ cm.\par
\tithelo Η θέση κάθε ευθείας ως προς τον κύκλο.\par
\possindeo Συγκρίνουμε την απόσταση με την ακτίνα $\rho=5$.\par
\apantisi Η ευθεία με $3<5$ είναι τέμνουσα (δύο κοινά σημεία). Η ευθεία με $5=5$ είναι εφαπτομένη (ένα κοινό σημείο). Η ευθεία με $7>5$ είναι εξωτερική (κανένα).\par
\elegxos Όσο η ευθεία απομακρύνεται, τα κοινά σημεία μειώνονται: $2$, $1$, $0$.

\paradeigma{Εφαπτομένη σε σημείο του κύκλου}
\tixero Κύκλος $(\mathrm{Ο},\rho)$ και σημείο $\mathrm{Α}$ πάνω του.\par
\tithelo Σχεδίαση της εφαπτομένης στο $\mathrm{Α}$.\par
\possindeo Φέρνουμε την ακτίνα $\mathrm{ΟΑ}$. Στο $\mathrm{Α}$ σχεδιάζουμε την ευθεία $\varepsilon$ που είναι κάθετη στην $\mathrm{ΟΑ}$ (με γνώμονα).\par
\apantisi Η $\varepsilon$ είναι η εφαπτομένη στο $\mathrm{Α}$.\par
\elegxos Η απόσταση του $\mathrm{Ο}$ από την $\varepsilon$ είναι $\mathrm{ΟΑ}=\rho$, άρα η $\varepsilon$ εφάπτεται.

\paradeigma{Κύκλος που εφάπτεται σε ευθεία}
\tixero Ευθεία $\varepsilon$ και σημείο $\mathrm{Α}$ πάνω της.\par
\tithelo Κύκλος που εφάπτεται στην $\varepsilon$ στο $\mathrm{Α}$.\par
\possindeo Φέρνουμε την κάθετη στην $\varepsilon$ στο $\mathrm{Α}$. Παίρνουμε πάνω της ένα οποιοδήποτε σημείο $\mathrm{Κ}$ και γράφουμε τον κύκλο $(\mathrm{Κ},\mathrm{ΚΑ})$.\par
\apantisi Ο κύκλος $(\mathrm{Κ},\mathrm{ΚΑ})$.\par
\elegxos Η απόσταση του $\mathrm{Κ}$ από την $\varepsilon$ είναι $\mathrm{ΚΑ}$, ίση με την ακτίνα. Υπάρχουν άπειροι τέτοιοι κύκλοι, ανάλογα με τη θέση του $\mathrm{Κ}$.

\paradeigma{Εφαπτόμενες στα άκρα χορδής}
\tixero Κύκλος $(\mathrm{Ο},\rho)$ και χορδή $\mathrm{ΑΒ}$.\par
\tithelo Οι εφαπτόμενες στα $\mathrm{Α}$ και $\mathrm{Β}$.\par
\possindeo Σχεδιάζουμε τις ακτίνες $\mathrm{ΟΑ}$ και $\mathrm{ΟΒ}$. Στο $\mathrm{Α}$ φέρνουμε κάθετη στην $\mathrm{ΟΑ}$ και στο $\mathrm{Β}$ κάθετη στην $\mathrm{ΟΒ}$. Αυτές είναι οι εφαπτόμενες και τέμνονται σε ένα σημείο $\mathrm{Μ}$.\par
\apantisi Οι εφαπτόμενες τέμνονται στο $\mathrm{Μ}$ και τα $\mathrm{ΜΑ}$, $\mathrm{ΜΒ}$ είναι τα εφαπτόμενα τμήματα.\par
\elegxos Το $\mathrm{Μ}$ είναι έξω από τον κύκλο.

\begin{center}
\begin{tikzpicture}[scale=0.9]
\draw[thick] (0,0) circle (1.6);
\fill (0,0) circle (1.6pt) node[above] {\small $\mathrm{Ο}$};
\fill (-1.028,-1.226) circle (1.6pt) node[left] {$\mathrm{Α}$};
\fill (1.028,-1.226) circle (1.6pt) node[right] {$\mathrm{Β}$};
\fill (0,-2.09) circle (1.6pt) node[below] {$\mathrm{Μ}$};
\draw (-1.028,-1.226)--(1.028,-1.226);
\draw[gblue,thick] (-1.75,-0.02)--(0,-2.09)--(1.75,-0.02);
\draw[dashed] (0,0)--(-1.028,-1.226) (0,0)--(1.028,-1.226);
\end{tikzpicture}
\end{center}

\tmima{Ασκήσεις}

\zoni{Το κατάλαβα}{τέμνουσα, εφαπτομένη, εξωτερική}

\ex Για κύκλο με ακτίνα $\rho=6$ cm και ευθείες με απόσταση από το κέντρο $d=4$ cm, $d=6$ cm και $d=9$ cm, πες αν η ευθεία είναι τέμνουσα, εφαπτομένη ή εξωτερική και πόσα κοινά σημεία έχει με τον κύκλο.

\ex Βρες πόσα κοινά σημεία έχουν ο κύκλος και η ευθεία όταν: α) $\rho=8$, $d=3$ β) $\rho=4$, $d=4$ γ) $\rho=3$, $d=6$ δ) $\rho=2{,}5$, $d=2{,}5$ ε) $\rho=10$, $d=12$.

\ex Σωστό ή λάθος; α) Η εφαπτομένη έχει ένα μόνο κοινό σημείο με τον κύκλο. β) Η εφαπτομένη είναι κάθετη στην ακτίνα του σημείου επαφής. γ) Αν μια ευθεία είναι εξωτερική, η απόστασή της από το κέντρο είναι μικρότερη από την ακτίνα. δ) Από ένα σημείο του κύκλου περνά μία μόνο εφαπτομένη. ε) Μια ευθεία μπορεί να έχει τρία κοινά σημεία με έναν κύκλο. στ) Κάθε ευθεία που περνά από το κέντρο είναι τέμνουσα.

\ex Σχεδίασε κύκλο $(\mathrm{Ο},4\text{ cm})$ και τρεις ευθείες που απέχουν από το $\mathrm{Ο}$ $3$ cm, $4$ cm και $5$ cm. Πες τη θέση της καθεμιάς.

\ex Σχεδίασε έναν κύκλο, πάρε ένα σημείο $\mathrm{Α}$ πάνω του και σχεδίασε την εφαπτομένη στο $\mathrm{Α}$.

\ex Κύκλος $(\mathrm{Ο},3\text{ cm})$ και ευθεία $\varepsilon$ με απόσταση $3$ cm από το $\mathrm{Ο}$. Πόσα κοινά σημεία έχουν; Πώς πρέπει να αλλάξει η απόσταση για να έχουν δύο κοινά σημεία;

\zoni{Το χρησιμοποιώ}{κήπος, τροχοί, μπάλες}

\ex \textbf{Το ποτιστικό.} Το ποτιστικό της θείας Μαρίας βρέχει σε ακτίνα $4$ m. Τρία ευθύγραμμα μονοπάτια απέχουν από αυτό $2{,}5$ m, $4$ m και $6$ m. Για κάθε μονοπάτι πες πόσα σημεία του βρέχονται.

\ex \textbf{Ο τροχός.} Ο Ηλίας κινεί το ποδήλατό του σε ευθύγραμμο δρόμο. Ο τροχός έχει ακτίνα $35$ cm. Πόσα σημεία του τροχού ακουμπούν στον δρόμο και πόσο ψηλά από το έδαφος βρίσκεται ο άξονας (το κέντρο) του τροχού;

\ex \textbf{Η μπάλα.} Ο Νίκος ακουμπά μια μπάλα του μπάσκετ ακτίνας $12$ cm στη γωνία δαπέδου και τοίχου. Πόσο απέχει το κέντρο της από τον τοίχο και πόσο από το δάπεδο; Τι ισχύει για τον τοίχο και το δάπεδο ως προς τον κύκλο που σχηματίζει η μπάλα στο πλάι;

\ex \textbf{Η λίμνη.} Στο πάρκο υπάρχει κυκλική λίμνη ακτίνας $50$ m. Ένας δρόμος περνά σε απόσταση $38$ m από το κέντρο της. Είναι εφικτός ο δρόμος; Πόση πρέπει να είναι η απόσταση ώστε ο δρόμος να περνά ακριβώς στην άκρη της λίμνης, χωρίς να τη διασχίζει;

\ex \textbf{Το σήμα.} Ο δρομολογητής Wi-Fi της Ρήνης έχει εμβέλεια $25$ m. Ο δρόμος απέναντι απέχει $30$ m από αυτόν. Φτάνει το σήμα στον δρόμο; Αν ο δρόμος ήταν σε απόσταση $25$ m, σε πόσα σημεία θα έφτανε; Και αν ήταν σε απόσταση $20$ m;

\zoni{Το σκέφτομαι}{κατασκευές και συμπεράσματα}

\ex Σε κύκλο $(\mathrm{Ο},\rho)$ φέρε δύο κάθετες διαμέτρους $\mathrm{ΑΒ}$ και $\mathrm{ΓΔ}$ και σχεδίασε τις εφαπτόμενες στα $\mathrm{Α}$, $\mathrm{Β}$, $\mathrm{Γ}$, $\mathrm{Δ}$. Τι σχήμα σχηματίζουν οι τέσσερις εφαπτόμενες και πόση είναι η πλευρά του;

\ex Σχεδίασε κύκλο $(\mathrm{Ο},4{,}5\text{ cm})$ και μια διάμετρο $\mathrm{ΑΒ}$. Σχεδίασε τις εφαπτόμενες του κύκλου που είναι παράλληλες προς την $\mathrm{ΑΒ}$. Πόσες είναι;

\ex Οι ευθείες $\varepsilon_1$ και $\varepsilon_2$ είναι εφαπτόμενες ενός κύκλου στα σημεία $\mathrm{Α}$ και $\mathrm{Β}$. Περίγραψε πώς μπορούμε να βρούμε το κέντρο του κύκλου.

\ex Πόσες εφαπτόμενες του κύκλου περνούν από ένα σημείο: α) πάνω στον κύκλο; β) έξω από τον κύκλο; γ) μέσα στον κύκλο; Δικαιολόγησε.

\zoni{Για το σπίτι}{ρίξε μια ματιά}

\ex Πάρε ένα κυκλικό καπάκι και έναν χάρακα. Τοποθέτησε τον χάρακα ώστε να εφάπτεται στο καπάκι. Μέτρησε την απόσταση του κέντρου από την άκρη του χάρακα. Τι βρίσκεις;

\ex Σχεδίασε ένα ποδήλατο: δύο ίσους κύκλους που εφάπτονται στην ίδια ευθεία (τον δρόμο). Πού βρίσκονται τα κέντρα τους ως προς τον δρόμο;
